\chapter{Stroke}

\section*{Overview}
A stroke occurs when blood flow to part of the brain is interrupted or reduced, preventing brain tissue from getting oxygen and nutrients. Brain cells begin to die in minutes. Stroke is a medical emergency. Prompt treatment minimizes brain damage and complications. There are two main types: ischemic stroke (blood clot) and hemorrhagic stroke (bleeding). TIA (transient ischemic attack) is a warning stroke.

\section*{Symptoms}
Stroke symptoms appear suddenly and may include:

\begin{itemize}
\item \textbf{Face drooping:} One side of the face droops or feels numb. Ask the person to smile.
\item \textbf{Arm weakness:} One arm is weak or numb. Ask the person to raise both arms.
\item \textbf{Speech difficulty:} Speech is slurred or strange. Ask the person to repeat a simple sentence.
\item \textbf{Time to call emergency services:} If anyone shows these symptoms, even if the symptoms go away, call emergency services immediately.
\end{itemize}

\textbf{Additional symptoms may include:}
\begin{itemize}
\item Numbness or weakness of the face, arm, or leg, especially on one side of the body
\item Confusion, trouble speaking, or understanding speech
\item Trouble seeing in one or both eyes
\item Trouble walking, dizziness, loss of balance or coordination
\item Severe headache with no known cause
\item Loss of consciousness
\end{itemize}

\section*{Causes}
\textbf{Ischemic stroke (87\% of strokes):}

Caused by blockage of blood vessels supplying the brain.

\begin{itemize}
\item \textbf{Thrombotic stroke:} Blood clot forms in an artery supplying blood to the brain, often due to atherosclerosis (plaque buildup)
\item \textbf{Embolic stroke:} Blood clot or debris forms away from the brain and travels through the bloodstream to lodge in narrower brain arteries
\end{itemize}

\textbf{Hemorrhagic stroke:}

Caused by bleeding in or around the brain.

\begin{itemize}
\item \textbf{Intracerebral hemorrhage:} Blood vessel bursts and leaks into surrounding brain tissue, often due to hypertension
\item \textbf{Subarachnoid hemorrhage:} Bleeding between the brain and surrounding membrane, often from aneurysm rupture
\end{itemize}

\textbf{Transient ischemic attack (TIA):} Temporary blockage causing stroke-like symptoms that resolve within 24 hours (usually minutes). Warning sign of future stroke.

\textbf{Risk factors:}
\begin{itemize}
\item High blood pressure (hypertension) - most important modifiable risk factor
\item Atrial fibrillation and other heart conditions
\item High cholesterol
\item Diabetes
\item Smoking
\item Obesity and physical inactivity
\item Family history
\item Age over 55 (risk doubles each decade after 55)
\item Previous stroke or TIA
\end{itemize}

\section*{Diagnosis}
Emergency diagnosis includes:

\textbf{Neurological examination:}
\begin{itemize}
\item Checking alertness, vision, arm and leg strength, coordination, and speech
\item \textbf{FAST assessment:} Face, Arms, Speech, Time
\item \textbf{NIH Stroke Scale:} Standardized assessment tool
\end{itemize}

\textbf{Imaging tests:}
\begin{itemize}
\item \textbf{CT scan:} Quickly identifies ischemic vs. hemorrhagic stroke; first-line imaging
\item \textbf{MRI:} Provides detailed images of brain tissue; more sensitive for early ischemia
\item \textbf{CT angiography:} Shows blood vessels in the neck and brain
\item \textbf{MR angiography:} MRI version of angiography
\item \textbf{CT perfusion:} Assesses blood flow to brain tissue
\end{itemize}

\textbf{Ultrasound:}
\begin{itemize}
\item \textbf{Carotid ultrasound:} Checks for plaque in neck arteries
\item \textbf{Echocardiogram:} Examines heart for clots or source of emboli
\end{itemize}

\textbf{Blood tests:} Blood sugar, clotting ability (PT/INR, aPTT), complete blood count, electrolytes, and other measurements.

\section*{Treatment}
Treatment depends on stroke type, severity, and timing:

\textbf{Ischemic stroke:}

\begin{itemize}
\item \textbf{tPA (tissue plasminogen activator):} Clot-busting drug given within 3-4.5 hours of symptom onset. Can improve outcomes but increases bleeding risk.
\item \textbf{Mechanical thrombectomy:} Physical removal of clot via catheter, effective up to 24 hours in select patients with large vessel occlusion
\item Antiplatelet drugs (aspirin, clopidogrel)
\item Anticoagulants (heparin, warfarin, direct oral anticoagulants) for cardioembolic strokes
\end{itemize}

\textbf{Hemorrhagic stroke:}

\begin{itemize}
\item \textbf{Blood pressure control:} Careful management to prevent rebleeding
\item \textbf{Surgical interventions:}
\begin{itemize}
\item Clipping of aneurysms
\item Coiling procedures (endovascular)
\item Removal of accumulated blood (evacuation)
\item Decompressive craniectomy
\item AVM (arteriovenous malformation) repair
\end{itemize}
\item Reversal of blood thinners if applicable
\item Management of intracranial pressure
\end{itemize}

\textbf{Rehabilitation:}
\begin{itemize}
\item Physical therapy for mobility and strength
\item Occupational therapy for daily activities
\item Speech therapy for communication and swallowing
\item Cognitive rehabilitation
\end{itemize}

\section*{Prevention}
Stroke prevention focuses on risk factor modification:

\begin{itemize}
\item \textbf{Control blood pressure:} Keep below 130/80 mmHg
\item \textbf{Manage cholesterol:} Statins and diet modification; target LDL <100 mg/dL (or <70 for high risk)
\item \textbf{Control diabetes:} Maintain target blood sugar levels (HbA1c <7\%)
\item \textbf{Quit smoking:} Smoking doubles stroke risk; cessation reduces risk over time
\item \textbf{Exercise regularly:} At least 150 minutes of moderate activity weekly
\item \textbf{Maintain healthy weight:} BMI between 18.5-24.9
\item \textbf{Limit alcohol:} Maximum one drink daily for women, two for men
\item \textbf{Treat atrial fibrillation:} Anticoagulation reduces stroke risk by 60-70\%
\item \textbf{Healthy diet:} Mediterranean or DASH diet, low sodium (<2300 mg/day)
\item \textbf{Medication adherence:} Take prescribed medications consistently
\end{itemize}

\section*{When to Seek Medical Care}
Call emergency services immediately if you or someone shows stroke symptoms. Remember \textbf{F.A.S.T}:

\begin{itemize}
\item \textbf{F}ace drooping
\item \textbf{A}rm weakness
\item \textbf{S}peech difficulty
\item \textbf{T}ime to call emergency services
\end{itemize}

\textbf{Do NOT:}
\begin{itemize}
\item Wait to see if symptoms improve
\item Drive yourself to the hospital
\item Give food or drink
\item Give aspirin without medical evaluation (could worsen hemorrhagic stroke)
\end{itemize}

Even if symptoms resolve (TIA), seek immediate evaluation as it's a warning sign of impending stroke. Time is brain - every minute counts.